\chapter{ANTECEDENTES}

% ===========================================================================
% Estado del arte y estudio de alternativas (IT-43). Aquí va el fundamento
% teórico y la comparación con lo que ya existe. Es lo que convierte el TFG en
% un trabajo de investigación y no solo "un programa que funciona".
% Sustituye los textos guía por tu redacción.
% ===========================================================================

\section{Sistemas de orientación y recomendación académica}

% IT-43. Comparar con lo que ya existe: orientación oficial, comparadores de
% carreras, chatbots de orientación de otras universidades. Probarlos de verdad,
% no describir su marketing. Conclusión: qué carencia cubre este TFG.

\section{Recuperación aumentada por generación (RAG)}

% Fundamento del enfoque: por qué RAG y no un LLM "a secas" ni un sistema de
% reglas. Concepto de recuperación + generación, y por qué reduce las respuestas
% inventadas en un dominio cerrado como la oferta académica.

\section{Modelos de lenguaje de código abierto}

% Alternativas de LLM de generación: Llama 3, Mistral, Phi-3 frente a modelos
% privados (GPT-4, Claude). Justificar el enfoque open-source en local
% (privacidad, coste, control). La decisión final se formaliza en el ADR-0005.

\section{Modelos de \textit{embeddings} para el español}

% Por qué hacen falta modelos multilingües o específicos de español, y no el
% MiniLM monolingüe inglés de los tutoriales genéricos. La elección concreta se
% decide experimentalmente (IT-21) y se formaliza en el ADR-0003.

\section{Bases de datos vectoriales}

% Alternativas: FAISS, ChromaDB, Qdrant (locales) frente a Pinecone (nube).
% Argumentar que, a la escala de este corpus, una solución local es suficiente.
% La elección se decide con datos (IT-24) y se formaliza en el ADR-0004.

\section{Segmentación de documentos (\textit{chunking})}

% Estrategias: tamaño fijo, por delimitadores, semántica. Es el fundamento del
% ADR-0001 (IT-11).
